\section*{Capítulo \ref{ch:b3:product} --- Medidas producto,
Fubini, cambio de variable}

\begin{solution}{exo:b3:product:1}
$\Delta$ es cerrado en $\intcc01^2$, luego de Borel, y $\mathcal
B(\intcc01^2)$ es la $\sigma$-álgebra producto
(la~\cref{prop:b3:product:sections}(b)). Iterando en un sentido:
\[
\int_{\intcc01}\Bigl(\int \mathbf 1_\Delta(x,y)
\,\dd\lambda(y)\Bigr)\dd\mu(x)
= \int \lambda(\{x\})\,\dd\mu(x) = 0 ;
\]
en el otro:
\[
\int_{\intcc01}\Bigl(\int\mathbf
1_\Delta(x,y)\,\dd\mu(x)\Bigr)\dd\lambda(y)
= \int \mu(\{y\})\,\dd\lambda(y) = \int 1\,\dd\lambda = 1 .
\]
La hipótesis que falla es la $\sigma$-finitud de la
\omterm{def:b3:measure:measure}{medida} de conteo $\mu$ sobre el
conjunto no numerable $\intcc01$: ninguna familia numerable de conjuntos
de $\mu$ finita lo recubre.
\end{solution}

\begin{solution}{exo:b3:product:2}
En $\intcc0A\times\intoo0{+\infty}$:
$\int_0^A\int_0^\infty\eu^{-xy}\abs{\sin x}\,\dd y\,\dd x =
\int_0^A\frac{\abs{\sin x}}x\dd x \leq A < \infty$ (Tonelli para el
valor absoluto): se aplica Fubini y, puesto que
$\int_0^\infty\eu^{-xy}\dd y = \frac1x$,
\[
\int_0^A\frac{\sin x}x\dd x
= \int_0^\infty\Bigl(\int_0^A\eu^{-xy}\sin x\,\dd
x\Bigr)\dd y
= \int_0^\infty \frac{1 - \eu^{-Ay}(\cos A + y\sin A)}{1 +
y^2}\,\dd y
\]
(la integral \omterm{def:b3:topology:interior}{interior}: $\operatorname{Im}\int_0^A\eu^{(\iu -
y)x}\dd x$, calculada directamente). Cuando $A \to \infty$, el término
correctivo está acotado por $\int_0^\infty\eu^{-Ay}\frac{1 +
y}{1 + y^2}\dd y \leq \frac32\int_0^\infty\eu^{-Ay}\dd y =
\frac3{2A} \to 0$; y el término principal es
$\int_0^\infty\frac{\dd y}{1+y^2} = \frac\pi2$. Por tanto,
$\int_0^\infty\frac{\sin x}x\dd x = \frac\pi2$ --- la integral de
Dirichlet por Fubini.
\end{solution}

\begin{solution}{exo:b3:product:3}
(a) Fórmula de las capas (la~\cref{prop:b3:product:layercake}): $\int
f\,\dd\mu = \int_0^\infty\mu(f > t)\,\dd t$, y $t \mapsto
\mu(f > t)$ es no creciente. En $[n-1, n]$: $\mu(f \geq n)
\leq \mu(f > t) \leq \mu(f > n - 1) \leq \mu(f \geq n - 1)$; sumando las
integrales sobre los intervalos unidad:
\[
\sum_{n\geq1}\mu(f \geq n) \leq \int f\,\dd\mu \leq
\sum_{n\geq1}\mu(f \geq n - 1) = \mu(f \geq 0) +
\sum_{n\geq1}\mu(f\geq n) \leq \mu(X) + \sum_{n\geq1}\mu(f\geq
n).
\]

(b) Aplíquese (a) a $\abs f$: la finitud de la integral y la de la serie
son equivalentes (el $\mu(X)$ adicional es finito).
\end{solution}

\begin{solution}{exo:b3:product:4}
Como $f(x,y) = \partial_y\bigl(\frac{y}{x^2+y^2}\bigr)$ para
$x \ne 0$:
\[
\int_0^1 f(x, y)\,\dd y = \frac{1}{x^2 + 1}
\ \Longrightarrow\
\int_0^1\!\!\int_0^1 f\,\dd y\,\dd x =
\int_0^1\frac{\dd x}{1 + x^2} = \frac\pi4 ;
\]
por la antisimetría $f(y,x) = -f(x,y)$, el otro orden da $-\frac\pi4$.
Valores absolutos: para $0 < y < x$,
\[
\int_0^x f(x,y)\,\dd y = \Bigl[\frac{y}{x^2 +
y^2}\Bigr]_0^x = \frac1{2x},
\quad\text{y } f \geq 0 \text{ allí, luego}\quad
\int_0^1\!\!\int_0^1\abs f \geq \int_0^1\frac{\dd x}{2x} =
+\infty .
\]
No hay contradicción con Fubini: su hipótesis $f \in L^1$ falla, y las
dos integrales iteradas son sencillamente dos números distintos.
\end{solution}

\begin{solution}{exo:b3:product:5}
(a) $(\mathbf 1_{\intcc01}*\mathbf 1_{\intcc01})(x) =
\lambda\bigl(\intcc01\cap\intcc{x-1}x\bigr)$: $0$ para $x
\notin \intcc02$, $x$ para $0 \leq x \leq 1$, $2 - x$ para $1
\leq x \leq 2$: la función triangular. Convolucionando de nuevo se
obtiene una joroba $\mathcal
C^1$ cuadrática a trozos sobre $\intcc03$ (la B-spline cuadrática): cada
convolución gana un grado de regularidad --- el principio de suavizado
que hay detrás de las aproximaciones de la identidad
del~\cref{ch:b3:lp}.

(b) Si $x \notin \overline{\operatorname{supp}f +
\operatorname{supp}g}$, hay una bola alrededor de $x$ disjunta del
conjunto suma; para $y \in \operatorname{supp}g$, $x - y
\notin\operatorname{supp}f$, de modo que el integrando se anula
idénticamente: $f * g = 0$ cerca de $x$.

(c) Para $x_n \to x$: $(f*g)(x_n) = \int
f(y)g(x_n - y)\,\dd y$; los integrandos convergen puntualmente
(\omterm{def:b3:topology:continuity}{continuidad} de $g$) y están
dominados por $\norm
g_\infty\,\abs f \in L^1$: la convergencia dominada da
$(f*g)(x_n) \to (f*g)(x)$.
\end{solution}

\begin{solution}{exo:b3:product:6}
(a) Por Fubini e inducción, cortando a lo largo de la última coordenada:
\[
\lambda_d(\Delta_d) = \int_0^1
\lambda_{d-1}\bigl((1 - t)\,\Delta_{d-1}\bigr)\,\dd t
= \lambda_{d-1}(\Delta_{d-1})\int_0^1(1 - t)^{d-1}\dd t
= \frac{\lambda_{d-1}(\Delta_{d-1})}{d},
\]
usando la regla de homotecia $\lambda_{d-1}(\rho A) =
\rho^{d-1}\lambda_{d-1}(A)$ (el~\cref{thm:b3:product:linearchange}); y
con $\lambda_1(\Delta_1) = 1$: volumen $\frac1{d!}$.

(b) $v_2 = \pi/\Gamma(2) = \pi$; $v_3 =
\pi^{3/2}/\Gamma(\frac52) = \pi^{3/2}/(\frac32\cdot\frac12
\sqrt\pi) = \frac{4\pi}3$. El elipsoide es $T(B(0,1))$ con
$T = \operatorname{diag}(a_1, \dots, a_d)$:
el~\cref{thm:b3:product:linearchange} da el volumen
$v_d\prod a_i$.
\end{solution}

\begin{solution}{exo:b3:product:7}
En $\R^2$, \omterm{ex:b3:product:polar}{coordenadas polares}
(el~\cref{ex:b3:product:polar}):
\[
\int_{B(0,1)}\frac{\dd x\dd y}{(x^2+y^2)^s}
= 2\pi\int_0^1 r^{1 - 2s}\,\dd r,
\qquad
\int_{\R^2\setminus B(0,1)} = 2\pi\int_1^\infty r^{1-2s}\dd r:
\]
finita si y solo si $1 - 2s > -1$ ($s < 1$), y respectivamente
$1 - 2s < -1$ ($s >
1$). En $\R^d$, evítense las coordenadas esféricas con la fórmula de las
capas: $\lambda_d(\{\norm x^{-2s} > t\}\cap B(0,1)) =
\lambda_d(B(0, \min(1, t^{-1/2s}))) = v_d\min(1, t^{-d/2s})$, y
$\int_0^\infty v_d\min(1, t^{-d/(2s)})\dd t < \infty$ si y solo si
$\frac d{2s} > 1$, es decir, $s < \frac d2$; la integral exterior
converge si y solo si $s > \frac d2$ (mismo cálculo en la región
complementaria).
\end{solution}

\begin{solution}{exo:b3:product:8}
Por Tonelli (integrandos positivos) y el cambio de variable
$(x, y) = \Phi(u, v) = (uv,\ u(1-v))$, a $\mathcal C^1$
difeomorfismo de $\intoo0\infty\times\intoo01$ sobre el cuadrante
abierto con
\[
\det D\Phi = \det\begin{pmatrix} v & u\\ 1 - v & -u
\end{pmatrix} = -uv - u(1 - v) = -u,
\qquad \abs{\det} = u :
\]
\[
\Gamma(p)\Gamma(q) =
\iint x^{p-1}y^{q-1}\eu^{-x-y}\dd x\,\dd y
= \iint (uv)^{p-1}\bigl(u(1{-}v)\bigr)^{q-1}\eu^{-u}\,u\,
\dd u\,\dd v
= \Gamma(p + q)\,B(p, q).
\]
Sustituyendo $t = \sin^2\theta$ en $B(p,q)$ se obtiene
$2\int_0^{\pi/2}\sin^{2p-1}\theta\cos^{2q-1}\theta\,\dd\theta =
B(p, q)$. Wallis: $W_n =
\int_0^{\pi/2}\sin^n\theta\,\dd\theta = \frac12B\bigl(\frac{n +
1}2, \frac12\bigr) = \frac{\Gamma(\frac{n+1}2)\sqrt\pi}
{2\,\Gamma(\frac n2 + 1)}$ --- por ejemplo, $W_{2n} =
\frac\pi2\cdot\frac{(2n)!}{4^n(n!)^2}$ usando
$\Gamma(n + \frac12) = \frac{(2n)!}{4^nn!}\sqrt\pi$.
\end{solution}

\begin{solution}{exo:b3:product:9}
Indicadores: $\int\mathbf 1_B\,\dd(T_*\mu) = T_*\mu(B) =
\mu(T^{-1}B) = \int\mathbf 1_B\circ T\,\dd\mu$; la linealidad lo
extiende a $g$ simple y la convergencia monótona a $g \geq 0$ --- la
fórmula de transferencia. Es puramente formal: reexpresa integrales
contra $T_*\mu$, pero no dice nada sobre \emph{qué} es $T_*\mu$. El
contenido del~\cref{thm:b3:product:linearchange} y
del~\cref{thm:b3:product:changeofvar} es la identificación
\[
\Phi_*\bigl(\lambda_d\restriction_U\bigr) =
\abs{\det D\Phi^{-1}}\,\lambda_d\restriction_V
\quad\text{(una medida de densidad)},
\]
es decir, un cálculo de la \omterm{def:b3:measure:measure}{medida} imagen de la
\omterm{def:b3:measure:lebesgueouter}{medida de Lebesgue} --- siendo el
ingrediente analítico la geometría diferencial de $\Phi$, y no el
formalismo de teoría de la \omterm{def:b3:measure:measure}{medida}.
\end{solution}

\begin{solution}{exo:b3:product:10}
Por Tonelli la gaussiana se factoriza, de modo que, con $G_1 =
\int_\R\eu^{-s^2}\dd s = \sqrt\pi$ y $\int_\R
s^2\eu^{-s^2}\dd s = \frac{\sqrt\pi}2$ (intégrese por partes):
\[
\int_{\R^d}x_1^2\,\eu^{-\norm x^2}\dd x =
\frac{\sqrt\pi}2\;\pi^{(d-1)/2} = \frac{\pi^{d/2}}2,
\qquad
\int_{\R^d}\norm x^2\eu^{-\norm x^2}\dd x =
d\cdot\frac{\pi^{d/2}}2
\]
(por simetría, $\norm x^2 = \sum_ix_i^2$ aporta $d$ términos iguales ---
la comprobación de coherencia). Para la
\omterm{def:b3:measure:measure}{medida} normalizada
$\pi^{-d/2}\eu^{-\norm x^2}\dd x$, el segundo momento es
$\frac d2$.
\end{solution}

\begin{solution}{exo:b3:product:11}
(a) $H = \Phi^{-1}(\intoo0\infty)$ para $\Phi(x, y) = f(x) -
y$ intersecado con $\{y > 0\}$: es
\omterm{def:b3:lebesgue:measurable}{medible}, pues $(x, y)
\mapsto f(x)$ y $(x,y)\mapsto y$ son
\omterm{def:b3:lebesgue:measurable}{medibles} sobre el producto
(composiciones con las proyecciones). La $x$-sección de $H$ es
$\intoo0{f(x)}$, de \omterm{def:b3:measure:measure}{medida} $f(x)$:
Tonelli integra las secciones,
\[
\lambda_{d+1}(H) = \int_{\R^d}\lambda_1\bigl(\intoo0{f(x)}
\bigr)\,\dd x = \int_{\R^d}f\,\dd\lambda_d .
\]

(b) La gráfica es $\{(x,y) : y \geq f(x)\} \cap \{y \leq
f(x)\}$, \omterm{def:b3:lebesgue:measurable}{medible}; sus $x$-secciones
son puntos, de \omterm{def:b3:measure:measure}{medida} $0$: Tonelli da
$\lambda_{d+1}(\text{grafo}) =
\int 0 = 0$.

(c) $S^{d-1}$ es la unión de las dos gráficas $y =
\pm\sqrt{1 - \abs{x'}^2}$ sobre la bola unidad de $\R^{d-1}$ (separando
la última coordenada): unión de dos conjuntos nulos por (b), luego nula.
\end{solution}

\begin{solution}{exo:b3:product:12}
En $\intoo01^2$, $\frac1{1 - xy} = \sum_{n\geq0}(xy)^n$ con términos no
negativos: Tonelli permite integrar término a término,
\[
\iint\frac{\dd x\,\dd y}{1 - xy}
= \sum_{n\geq0}\Bigl(\int_0^1x^n\dd x\Bigr)
\Bigl(\int_0^1y^n\dd y\Bigr)
= \sum_{n\geq0}\frac1{(n+1)^2} = \zeta(2) .
\]
Para el caso alternado, $\frac1{1 + xy} =
\sum_n(-1)^n(xy)^n$ no es una serie de términos positivos; pero la
integral de la serie de los \emph{valores absolutos} es $\zeta(2) <
\infty$, de modo que Fubini (con la
\omterm{def:b3:lebesgue:l1}{integrabilidad} ya establecida) se aplica:
$\iint\frac{\dd x\dd y}{1 + xy} =
\sum_n\frac{(-1)^n}{(n+1)^2}$. La identidad de series:
\[
\sum_{n\geq1}\frac{(-1)^{n-1}}{n^2}
= \sum_{n\geq1}\frac1{n^2} - 2\sum_{k\geq1}\frac1{(2k)^2}
= \zeta(2) - \frac{\zeta(2)}2 = \frac{\zeta(2)}2 .
\]
Con $\zeta(2) = \frac{\pi^2}6$ (el~\cref{pb:b3:spectral:1}): las
integrales valen $\frac{\pi^2}6$ y $\frac{\pi^2}{12}$. La positividad de
Tonelli lo era todo en el primer cálculo --- no hacía falta comprobar
ninguna \omterm{def:b3:lebesgue:l1}{integrabilidad} antes de
intercambiar; en el segundo, la positividad de la serie de valores
absolutos es lo que \emph{certifica} la
\omterm{def:b3:lebesgue:l1}{integrabilidad} para que Fubini pueda actuar
sobre la serie con signo.
\end{solution}

\begin{solution}{pb:b3:product:1}
\textbf{1.} Con $x = t + \sqrt t\,u$ ($\dd x = \sqrt
t\,\dd u$; $x$ recorre $\intoo0\infty$ cuando $u$ recorre
$\intoo{-\sqrt t}\infty$):
\[
\Gamma(t{+}1) = \int_0^\infty x^t\eu^{-x}\dd x
= t^t\eu^{-t}\sqrt t\int_{-\sqrt t}^{\infty}
\exp\Bigl(t\ln\Bigl(1 + \frac u{\sqrt t}\Bigr) - \sqrt
t\,u\Bigr)\dd u,
\]
puesto que $x^t = t^t\exp\bigl(t\ln(1 + u/\sqrt t)\bigr)$ y
$\eu^{-x} = \eu^{-t}\eu^{-\sqrt tu}$.

\textbf{2.} Para $u$ fijo y $t \to \infty$: $t\ln(1 +
u/\sqrt t) - \sqrt tu = t\bigl(\frac u{\sqrt t} -
\frac{u^2}{2t} + o(\frac1t)\bigr) - \sqrt tu = -\frac{u^2}2 +
o(1)$: $g_t(u) \to \eu^{-u^2/2}$.

\textbf{3.} Póngase $\psi_1(h) = \varphi(h) + \frac{h^2}4$ en
$\intoc{-1}1$: $\psi_1(0) = 0$ y $\psi_1'(h) = \frac1{1+h} -
1 + \frac h2 = \frac{h(h-1)}{2(1+h)}$, que es $\geq 0$ en $\intoc{-1}0$
y $\leq 0$ en $\intcc01$: $\psi_1 \leq 0$, es decir,
$\varphi(h) \leq -h^2/4$ allí. Póngase $\psi_2(h) =
\varphi(h) + ch$ en $\intco1\infty$, $c = 1 - \ln2$: $\psi_2(1)
= \ln2 - 1 + c = 0$ y $\psi_2'(h) = c - \frac h{1+h} \leq c -
\frac12 < 0$: $\varphi(h) \leq -ch$ para $h \geq 1$. Ahora, para $t
\geq 1$: si $\abs u \leq \sqrt t$, $g_t(u) =
\eu^{t\varphi(u/\sqrt t)} \leq \eu^{-t(u/\sqrt t)^2/4} =
\eu^{-u^2/4}$; si $u \geq \sqrt t$, entonces $t\,\varphi(u/\sqrt t) \leq
-ct\cdot\frac u{\sqrt t} = -c\sqrt t\,u \leq -cu$ (pues $t \geq
1$), luego $g_t(u) \leq \eu^{-cu}$. Por tanto, $g_t \leq \eu^{-u^2/4} +
\eu^{-cu}\mathbf 1_{u>0}$, \omterm{def:b3:lebesgue:l1}{integrable} e
independiente de $t \geq
1$.

\textbf{4.} Convergencia dominada: $\int_\R g_t(u)\dd u \to
\int_\R\eu^{-u^2/2}\dd u = \sqrt{2\pi}$ (el~\cref{ex:b3:product:polar}
más el cambio de escala $u\mapsto\sqrt2\,u$). Con la pregunta 1:
\[
\Gamma(t + 1) \sim \sqrt{2\pi t}\;\Bigl(\frac
t\eu\Bigr)^t,\qquad
n! \sim \sqrt{2\pi n}\,\Bigl(\frac n\eu\Bigr)^n .
\]

\textbf{5.} Del~\cref{exo:b3:product:8}, $W_{2n} =
\frac\pi2\,\frac{(2n)!}{4^n(n!)^2} =
\frac\pi2\,4^{-n}\binom{2n}n$. Stirling:
\[
\binom{2n}{n} = \frac{(2n)!}{(n!)^2}
\sim \frac{\sqrt{4\pi n}\,(2n/\eu)^{2n}}
{2\pi n\,(n/\eu)^{2n}} = \frac{4^n}{\sqrt{\pi n}} .
\]
Entonces $W_{2n} \sim \frac12\sqrt{\pi/n}$, coherente con la recursión
$W_n = \frac{n-1}nW_{n-2}$ (que obliga a $W_n \sim
W_{n-2}$ y, con $W_nW_{n-1}\cdot n = \frac\pi2$ --- la relación clásica
de Wallis ---, da $W_n \sim
\sqrt{\pi/(2n)}$; las dos asintóticas concuerdan).

\textbf{6.} $\Gamma(\frac d2 + 1) \sim \sqrt{2\pi\frac
d2}\,(\frac d{2\eu})^{d/2}$, luego
\[
v_d = \frac{\pi^{d/2}}{\Gamma(\frac d2 + 1)}
\sim \frac{1}{\sqrt{\pi d}}\Bigl(\frac{2\pi\eu}
d\Bigr)^{d/2} \longrightarrow 0
\]
supergeométricamente (para $d > 2\pi\eu \approx 17$, cada factor es
$< 1$ y decreciente). Numéricamente $v_1 = 2$, $v_2 \approx
3.14$, $v_3 \approx 4.19$, $v_4 \approx 4.93$, $v_5 \approx
5.26$, $v_6 \approx 5.17$: el máximo está en $d = 5$.

\textbf{7.} Con $k = \frac n2 + \frac{s\sqrt n}2$ (entero, $n$ par, $s$
fijo): tómense logaritmos en $2^{-n}\binom nk =
2^{-n}\frac{n!}{k!(n-k)!}$ y aplíquese Stirling a los tres factoriales.
Escribiendo $k = \frac n2(1 + \varepsilon)$, $n - k =
\frac n2(1 - \varepsilon)$ con $\varepsilon = s/\sqrt n$:
\[
\ln\Bigl(2^{-n}\binom nk\Bigr)
= -\frac n2\bigl[(1{+}\varepsilon)\ln(1{+}\varepsilon) +
(1{-}\varepsilon)\ln(1{-}\varepsilon)\bigr]
+ \frac12\ln\frac{2}{\pi n(1 - \varepsilon^2)} + o(1),
\]
y el corchete es $\varepsilon^2 + O(\varepsilon^4) =
\frac{s^2}n + O(n^{-2})$: la expresión tiende a $-\frac{s^2}2 +
\frac12\ln\frac2{\pi n}$ salvo $o(1)$, es decir,
\[
2^{-n}\binom nk \sim \sqrt{\frac{2}{\pi
n}}\;\eu^{-s^2/2} :
\]
el perfil gaussiano del lanzamiento de monedas, cuantificado --- la
forma local de de Moivre--Laplace, que se globalizará en
\cref{ch:b3:clt}.

\textbf{8.} (i) La convergencia dominada convierte el límite puntual de
la pregunta 2 en convergencia de las integrales, usando el dominador de
la pregunta 3. (ii) La integral gaussiana evalúa el límite
$\int\eu^{-u^2/2} = \sqrt{2\pi}$ --- la constante $\sqrt{2\pi}$ de
Stirling \emph{es} la integral gaussiana. (iii) La sustitución
$x = t + \sqrt tu$ es un cambio de variable afín: invariancia por
traslaciones y regla de escala de la
\omterm{def:b3:measure:lebesgueouter}{medida de Lebesgue}
(el~\cref{thm:b3:product:linearchange} en dimensión $1$).

\textbf{9.} Desarróllense ambos términos:
\[
d_n - d_{n+1} = \ln\frac{n!}{(n+1)!} + \Bigl(n +
\frac32\Bigr)\ln(n+1) - \Bigl(n + \frac12\Bigr)\ln n - 1
= \Bigl(n + \frac12\Bigr)\ln\frac{n+1}n - 1,
\]
combinándose los términos $-\ln(n+1)$ y $(n + \frac32)\ln(n+1)$ en
$(n + \frac12)\ln(n+1)$.

\textbf{10.} Para $t = \frac1{2n+1}$: $\frac{1+t}{1-t} =
\frac{2n+2}{2n} = \frac{n+1}n$, y $n + \frac12 =
\frac1{2t}$; la serie impar $\ln\frac{1+t}{1-t} =
2\sum_{k\geq0}\frac{t^{2k+1}}{2k+1}$ da
\[
\Bigl(n + \frac12\Bigr)\ln\frac{n+1}n
= \sum_{k\geq0}\frac{t^{2k}}{2k+1}
= 1 + \frac{t^2}3 + \frac{t^4}5 + \cdots
\]
Réstese $1$. Cota inferior: el primer término solo,
$\frac{t^2}3 = \frac1{3(2n+1)^2}$. Cota superior: bájense todos los
denominadores a $3$ y súmese la serie geométrica:
$\frac{t^2}{3(1 - t^2)} = \frac1{3((2n+1)^2 - 1)} =
\frac1{12n(n+1)} = \frac1{12n} - \frac1{12(n+1)}$.

\textbf{11.} Sumando la cota superior desde $n$ hasta $\infty$ (con
$d_m \to 0$): $d_n < \frac1{12n}$. Para la cota inferior:
$\frac1{12m+1} - \frac1{12(m+1)+1} =
\frac{12}{(12m+1)(12m+13)}$, y
\begin{align*}
\frac1{3(2m+1)^2} > \frac{12}{(12m+1)(12m+13)}
&\iff (12m+1)(12m+13) > 36(2m+1)^2 \\
&\iff 168m + 13 > 144m + 36,
\end{align*}
cierto para $m \geq 1$. Sumando este minorante telescópico:
$d_n > \frac1{12n+1}$. Exponenciando se obtiene el encaje clásico de
$n!$.

\textbf{12.} (a) Error relativo $= \eu^{d_n} - 1 <
\eu^{1/(12n)} - 1 < \frac{1.1}{12n}$ para $n$ grande; $<
10^{-6}$ en cuanto $12n \geq 1.1\cdot10^6$, y el $n \geq 83\,334$
indicado basta ($\frac1{12n} \leq 10^{-6}$ ya lo implica). (b)
$\log_{10}(100!) =
\frac12\log_{10}(200\pi) + 200 - 100\log_{10}\eu +
d_{100}\log_{10}\eu = 1.39906 + 200 - 43.42945 + 0.00036
\approx 157.96997$, luego $100! \approx 10^{0.96997} \cdot
10^{157} = 9.333\cdot10^{157}$; la ventana garantizada
$(\eu^{1/1201}, \eu^{1/1200})$ tiene una anchura relativa inferior a
$10^{-6}$ --- una fórmula «asintótica» que, en $n =
100$, es un instrumento de precisión.

\textbf{13.} Escríbase $\sin^n = \sin^{n-2} - \sin^{n-2}\cos^2$ e
intégrese el segundo término por partes ($u = \cos\theta$,
$\dd v = \sin^{n-2}\cos\theta\,\dd\theta$, $v =
\frac{\sin^{n-1}}{n-1}$):
\[
\int_0^{\pi/2}\sin^{n-2}\cos^2 =
\Bigl[\cos\theta\,\frac{\sin^{n-1}\theta}{n-1}\Bigr]_0^{\pi/2}
+ \frac1{n-1}\int_0^{\pi/2}\sin^n = \frac{W_n}{n-1} .
\]
Por tanto, $W_n = W_{n-2} - \frac{W_n}{n-1}$, es decir, $W_n =
\frac{n-1}nW_{n-2}$. De $W_0 = \frac\pi2$, $W_1 = 1$:
\[
W_{2n} = \frac{(2n-1)!!}{(2n)!!}\cdot\frac\pi2
= \frac\pi2\binom{2n}n4^{-n},
\qquad
W_{2n+1} = \frac{(2n)!!}{(2n+1)!!}
= \frac{4^n}{(2n+1)\binom{2n}n},
\]
convirtiendo los factoriales dobles mediante $(2n)!! = 2^nn!$ y
$(2n-1)!! = \frac{(2n)!}{2^nn!}$. Por último, $nW_nW_{n-1} =
(n-1)W_{n-1}W_{n-2}$ por la recursión: constante, igual a
$1\cdot W_1W_0 = \frac\pi2$.

\textbf{14.} $W_{2n+1} \leq W_{2n} \leq W_{2n-1}$ (monotonía puntual de
$\sin^n$) y $\frac{W_{2n-1}}{W_{2n+1}} =
\frac{2n+1}{2n} \to 1$ encajan $\frac{W_{2n}}{W_{2n+1}} \to
1$. Combinado con $W_{2n}W_{2n+1} = \frac{\pi}{2(2n+1)}$ (pregunta 13):
$W_{2n}^2 \sim \frac\pi{4n}$, luego $W_{2n} \sim
\frac12\sqrt{\frac\pi n}$ y $\binom{2n}n4^{-n} =
\frac2\pi W_{2n} \sim \frac1{\sqrt{\pi n}}$.

\textbf{15.} Las preguntas 9--10 nunca usaron el valor de la constante:
con $e_n = \ln n! - (n+\frac12)\ln n + n$, las diferencias
$e_n - e_{n+1}$ están en $\bigl(0, \frac1{12n} -
\frac1{12(n+1)}\bigr)$, de modo que $(e_n)$ decrece mientras $(e_n -
\frac1{12n})$ crece: sucesiones adyacentes, convergentes a un mismo
$\ell$. Por tanto, $n! \sim K n^{n+1/2}\eu^{-n}$, $K =
\eu^\ell$.

\textbf{16.} Sustituyendo el Stirling con constante desconocida en el
binomial central:
\[
\binom{2n}n4^{-n}\sqrt n \sim
\frac{K\,(2n)^{2n+1/2}\eu^{-2n}}{\bigl(K\,n^{n+1/2}
\eu^{-n}\bigr)^2}\,4^{-n}\sqrt n
= \frac{2^{2n}\sqrt{2}\,K\,n^{2n+1/2}}{K^2\,n^{2n+1}}
\,4^{-n}\sqrt n = \frac{\sqrt2}{K},
\]
y la pregunta 14 obliga a $\frac{\sqrt2}K = \frac1{\sqrt\pi}$:
$K = \sqrt{2\pi}$. Las Partes III--IV vuelven, pues, a demostrar
Stirling desde cero; e introduciéndolo en la identidad de la Parte I se
evalúa $\int_\R\eu^{-u^2/2}\dd u = \sqrt{2\pi}$ sin
\omterm{ex:b3:product:polar}{coordenadas polares}: Wallis y Gauss se
sostienen mutuamente.

\textbf{17.} $\frac{\binom{2n}{n+j}}{\binom{2n}n} =
\frac{(n!)^2}{(n+j)!\,(n-j)!} =
\prod_{i=1}^{j}\frac{n-i+1}{n+i}$ para $j \geq 0$ (y por simetría para
$j < 0$). Tomando logaritmos, con $1 \leq i
\leq j \leq K\sqrt n$:
\[
\ln\frac{n-i+1}{n+i} = \ln\Bigl(1 - \frac{i-1}n\Bigr) -
\ln\Bigl(1 + \frac in\Bigr) = -\frac{2i-1}{n} +
O\Bigl(\frac{i^2}{n^2}\Bigr),
\]
y $\sum_{i\leq j}(2i - 1) = j^2$, mientras que el error suma
$O(j^3/n^2) = O(n^{-1/2})$: uniformemente,
$\exp\bigl(-\frac{j^2}n + O(n^{-1/2})\bigr)$.

\textbf{18.} $\eu^{-n}\frac{n^n}{n!} \sim \eu^{-n}
\frac{n^n}{\sqrt{2\pi n}\,n^n\eu^{-n}} =
\frac1{\sqrt{2\pi n}}$. Una variable de Poisson de media $n$ tiene
desviación típica $\sqrt n$, y $\frac1{\sqrt{2\pi n}}$ es exactamente la
altura del pico gaussiano $\frac1{\sigma\sqrt{2\pi}}$: el teorema
central del límite local, anticipado en la moda.

\textbf{19.} Para $a \in \intoo01$, el lema de las pendientes
del~\cref{pb:b3:lebesgue:1} (su pregunta 14), aplicado a la convexa
$\log\Gamma$ alrededor de $n$, da $(n-1)^a \leq
\frac{\Gamma(n+a)}{\Gamma(n)} \leq n^a$: el cociente respecto de $n^a$
queda encajado por $(1 - \frac1n)^a \to 1$. Para $a = m + a'$
($m \in \N$, $a' \in \intco01$): $\Gamma(n+a) = (n + a - 1)
\cdots(n + a')\Gamma(n + a')$, y cada uno de los $m$ factores es
$n(1 + O(\frac1n))$: multiplíquense las estimaciones.

\textbf{20.} La recursión $v_d = \frac{2\pi}dv_{d-2}$ (de
$\Gamma(\frac d2 + 1) = \frac d2\Gamma(\frac d2)$) da, a partir de
$v_1 = 2$, $v_2 = \pi$:
\[
v_3 = \frac{4\pi}3,\quad v_4 = \frac{\pi^2}2,\quad
v_5 = \frac{8\pi^2}{15},\quad v_6 = \frac{\pi^3}6,\quad
v_7 = \frac{16\pi^3}{105}.
\]
El cociente $\frac{2\pi}d$ supera $1$ exactamente para $d \leq 6$, de
modo que cada paridad crece y después decrece; numéricamente $v_4
\approx 4.93$, $v_5 \approx 5.26$, $v_6 \approx 5.17$: el máximo global
es $d = 5$. Función generatriz: $v_{2k} =
\frac{\pi^k}{k!}$, luego $\sum_kv_{2k}x^{2k} =
\eu^{\pi x^2}$ --- todos los volúmenes de bolas de dimensión par
enrollados en una sola exponencial, y una estimación inmediata de
decaimiento supergeométrico para $v_d$.

\textbf{21.} Stirling en el numerador y en el denominador, con $k =
\alpha n$:
\[
\binom n{\alpha n} \sim
\frac{\sqrt{2\pi n}\,n^n}
{\sqrt{2\pi\alpha n}\,(\alpha n)^{\alpha n}\,
\sqrt{2\pi(1-\alpha)n}\,((1-\alpha)n)^{(1-\alpha)n}}
= \frac{\eu^{nH(\alpha)}}{\sqrt{2\pi\alpha(1-\alpha)n}},
\]
puesto que $n^n/(\alpha n)^{\alpha n}((1-\alpha)n)^{(1-\alpha)n} =
\alpha^{-\alpha n}(1-\alpha)^{-(1-\alpha)n} =
\eu^{nH(\alpha)}$ (las potencias de $n$ se cancelan: $\alpha n +
(1-\alpha)n = n$), y las $\eu^{-n}$ se cancelan igualmente. En
$\alpha = \frac12$: $H = \ln2$ y el prefactor es $\sqrt{2/(\pi n)}$ ---
de nuevo la pregunta 5. La concavidad estricta de $H$ (su derivada
segunda es $-\frac1{\alpha(1-\alpha)} < 0$) sitúa su máximo $\ln 2$ solo
en $\alpha = \frac12$: para $\alpha \neq \frac12$,
$\binom n{\alpha n}2^{-n} \approx
\eu^{-n(\ln2 - H(\alpha))}$ decae exponencialmente --- el motor
combinatorio que hay detrás de todo enunciado de concentración sobre
lanzamientos de moneda.

\textbf{22.} De $s_{d-1} = dv_d$ y la pregunta 20:
\[
s_0 = 2,\ \ s_1 = 2\pi,\ \ s_2 = 4\pi,\ \ s_3 = 2\pi^2,\ \
s_4 = \frac{8\pi^2}3,\ \ s_5 = \pi^3,\ \ s_6 =
\frac{16\pi^3}{15},
\]
numéricamente $2,\ 6.28,\ 12.57,\ 19.74,\ 26.32,\ 31.01,\
33.07$; y $s_7 = \frac{\pi^4}3 \approx 32.47 < s_6$: el máximo es la
$6$-esfera. La recursión $s_{d+1} =
(d+2)\,v_{d+2} = (d+2)\,\frac{2\pi}{d+2}\,v_d = 2\pi v_d =
\frac{2\pi}d\,s_{d-1}$ muestra la misma subida impulsada por
$\frac{2\pi}d$ y la misma caída supergeométrica que en los volúmenes:
pasada la dimensión siete, las esferas se encogen más deprisa que
cualquier sucesión geométrica.

\textbf{23.} La pregunta 11 dice exactamente $\frac1{12n+1} < d_n <
\frac1{12n}$, y
\[
\frac1{12n} - \frac1{12n+1} = \frac1{12n(12n+1)} =
O\Bigl(\frac1{n^2}\Bigr),
\]
luego $d_n = \frac1{12n} + O(\frac1{n^2})$ y $\eu^{d_n} = 1 +
\frac1{12n} + O(\frac1{n^2})$; multiplicando por
$\sqrt{2\pi n}(n/\eu)^n$ se obtiene la fórmula corregida. En $n =
10$: $\sqrt{20\pi}\,(10/\eu)^{10} = 7.92665 \times 453999.3
\approx 3\,598\,696$, bajo por $30\,104$ (error relativo
$8.3\cdot10^{-3}$); multiplicando por $1 + \frac1{120}$ se obtiene
$3\,628\,685$, bajo por $115$ (error relativo $3.2\cdot10^{-5}$). El
propio encaje sitúa $10!$ entre
$3\,598\,696\,\eu^{1/121} \approx 3\,628\,559$ y
$3\,598\,696\,\eu^{1/120} \approx 3\,628\,808$ --- la cota superior se
desvía en ocho unidades de la séptima cifra.

\textbf{24.} La sustitución $x = t + \sqrt t\,u$ de la Parte I, aplicada
a la integral truncada, da
\[
\int_0^{t} x^{t}\eu^{-x}\,\dd x
= t^{t}\eu^{-t}\sqrt t\int_{-\sqrt t}^{0} g_t(u)\,\dd u ,
\]
convirtiéndose el rango $0 \leq x \leq t$ en $-\sqrt t \leq u \leq
0$. El dominador de la pregunta 3 cubre también $g_t\mathbf 1_{u < 0}$,
de modo que la convergencia dominada da
\[
\int_{-\sqrt t}^{0}g_t(u)\,\dd u \longrightarrow
\int_{-\infty}^{0}\eu^{-u^2/2}\dd u = \frac{\sqrt{2\pi}}2 ,
\]
mientras que la pregunta 4 da $\Gamma(t+1) \sim
t^t\eu^{-t}\sqrt t\,\sqrt{2\pi}$. El cociente tiende a $\frac12$. En
términos probabilísticos: una variable aleatoria Gamma de parámetro de
forma grande pone asintóticamente la mitad de su masa a cada lado de su
moda --- la simetría del teorema central del límite, leída en una sola
sustitución.

\textbf{25.} Sea $\lambda = \frac{\alpha}{1-\alpha} \in
\intoc01$. Para $k \leq \lfloor\alpha n\rfloor \leq \alpha n$ se tiene
$\lambda^{k - \alpha n} \geq 1$, luego
\[
\sum_{k=0}^{\lfloor\alpha n\rfloor}\binom nk
\leq \lambda^{-\alpha n}\sum_{k=0}^{n}\binom nk\lambda^{k}
= \Bigl(\lambda^{-\alpha}(1 + \lambda)\Bigr)^{n},
\]
y con $\lambda = \frac\alpha{1-\alpha}$:
\[
\lambda^{-\alpha}(1+\lambda)
= \alpha^{-\alpha}(1-\alpha)^{\alpha}\cdot\frac1{1-\alpha}
= \alpha^{-\alpha}(1-\alpha)^{-(1-\alpha)}
= \eu^{H(\alpha)} .
\]
Optimalidad: minimizando $f(\lambda) = -\alpha\ln\lambda +
\ln(1+\lambda)$ sobre $\lambda > 0$, la ecuación $f'(\lambda)
= -\frac\alpha\lambda + \frac1{1+\lambda} = 0$ tiene la única solución
$\lambda = \frac\alpha{1-\alpha}$, que es un mínimo porque $f'' > 0$ ---
la elección del sesgo exponencial (Chernoff). Reconciliación: por la
pregunta 21, el término aislado $k =
\lfloor\alpha n\rfloor$ ya es del orden de
$\eu^{nH(\alpha)}/\sqrt{2\pi\alpha(1-\alpha)n}$, de modo que
\[
\frac{\eu^{nH(\alpha)}}{C\sqrt n} \leq
\sum_{k\leq\alpha n}\binom nk \leq \eu^{nH(\alpha)} :
\]
la tasa $H(\alpha)$ es exacta, y toda la suma cuesta a lo sumo un factor
$\sqrt n$ sobre su término mayor. Dividido por $2^n$, esto es la cota de
cola para la moneda equilibrada $\P(S_n \leq \alpha n) \leq
\eu^{-n(\ln2 - H(\alpha))}$ --- concentración de la
\omterm{def:b3:measure:measure}{medida} en una línea.
\end{solution}
